\subsection{Binary Access Patterns}

Binary workloads emphasize byte counts, fixed chunk sizes, and seek positions rather than lines.

\subsubsection{Reading Fixed-Size Chunks for Large Files}

\texttt{read(n)} in binary mode returns at most \texttt{n} bytes; empty \texttt{b""} signals EOF. Only one chunk resides in memory at a time.

%<<<PROGRAM:programs/bin01>>>
\subsubsection{Writing Byte Buffers to External Storage}

\texttt{write()} accepts \texttt{bytes} and \texttt{bytearray}; return value is bytes written.

\subsubsection{Random Access with \texttt{.seek()} and \texttt{.tell()}}

\texttt{tell()} reports the current byte offset; \texttt{seek()} repositions before read/write. Overwrites change bytes in place without list-like insertion semantics.

\begin{verbatim}
path = "seek_demo.bin"
with open(path, "wb") as f:
    f.write(b"0123456789")
with open(path, "r+b") as f:
    f.seek(5)
    print(f.read(3))
\end{verbatim}
